\chapter{Theoretische Grundlagen und Stand der Technik}
\section{Explorative Datenanalyse (EDA)}

Der Begriff der explorativen Datenanalyse geht auf \textcite{tukey1977eda} zurück, der damit einen Analyseansatz prägte, bei dem Daten zunächst offen und grafisch gestützt untersucht werden, um Muster und Auffälligkeiten zu erkennen und Hypothesen zu gewinnen, statt vorab festgelegte Annahmen konfirmatorisch zu prüfen. Heute wird die \gls{eda} als ein iterativer Analyseprozess verstanden, in dessen Verlauf ein Datensatz systematisch untersucht wird, um Strukturen zu erkennen, Anomalien wie Ausreißer aufzudecken sowie erste Hypothesen zu generieren und vorläufig z u bewerten \parencite{langer_towards_2019}. Im Rahmen dieser Methodik kommen primär deskriptive Statistiken und visuelle Werkzeuge zum Einsatz, um aussagekräftige Hinweise aus den Daten zu gewinnen \parencite{rahmany_comparing_2020}.

Ein zentrales Merkmal der \gls{eda} ist ihr iterativer Charakter. Sie stellt keinen einmaligen, abgeschlossenen Prozess dar, sondern umfasst mehrere aufeinander aufbauende Analyseschritte, bis ein hinreichendes Verständnis der Daten erreicht ist \parencite{rahmany_comparing_2020}. Da dieser Prozess nicht rein automatisiert abläuft, hängt die Qualität der Ergebnisse maßgeblich von der Fähigkeit des Analysten ab, die richtigen Fragen zu stellen und eine geeignete Analyseinteraktion durchzuführen, um die Daten sinnvoll zu interpretieren \parencite{langer_towards_2019}.

Zur strukturierten Durchführung beschreiben \textcite{rahmany_comparing_2020} sechs aufeinander aufbauende Schritte. Der Prozess beginnt mit der formalen Erfassung, bei der die vorhandenen Variablen nach Datentyp und Bedeutung erfasst und klassifiziert werden. Darauf aufbauend erfolgt eine univariate Analyse der einzelnen Variablen hinsichtlich zentraler Tendenz und Streuung. Im dritten Schritt werden bivariate und multivariate Zusammenhänge zwischen den Attributen untersucht, insbesondere Korrelationen und Kovarianzen. Darauf folgt die Identifikation auffälliger Muster wie fehlender oder inkonsistenter Werte. Im fünften Schritt werden Ausreißer erkannt, also Datenpunkte, die deutlich von der übrigen Verteilung abweichen und Analyseergebnisse verzerren können. Abschließend umfasst das Feature Engineering die Aufbereitung der Rohdaten für weiterführende Analysen.

Unabhängig von der genauen Prozessphase lassen sich die in der \gls{eda} eingesetzten Techniken grundsätzlich in zwei Kategorien unterteilen: deskriptive Statistiken zur quantitativen Beschreibung der Daten sowie grafische Methoden zur visuellen Exploration \parencite{rahmany_comparing_2020}.

\section{Large Language Models (LLMs)}
\label{sec:llm}

\glspl{llm} sind neuronale Sprachmodelle, die auf Basis der Transformer-Architektur trainiert werden \parencite{vaswani_attention_2017}. Die Wissensgrundlage dieser Modelle entsteht in der Vortrainingsphase, in der das Modell auf umfangreichen Textkorpora trainiert wird, um die Wahrscheinlichkeitsverteilung über Tokenfolgen zu erlernen und sukzessive den jeweils nächsten Token einer Sequenz vorherzusagen \parencite{brown_language_2020}. Durch diesen Prozess werden sprachliche Muster, faktisches Wissen und prozedurale Zusammenhänge implizit in den Modellparametern kodiert. Zur Beantwortung einer Anfrage greift das Modell nicht auf eine externe Wissensquelle zu, sondern ausschließlich auf das in den Gewichten gespeicherte Wissen \parencite{brown_language_2020}. Aus dieser rein parametrischen Wissensrepräsentation resultiert jedoch eine methodische Einschränkung. Da das Modell nicht zwischen erlerntem Vorwissen und den im Prompt bereitgestellten Daten unterscheidet, besteht die Gefahr, dass Ausgaben stärker durch statistisch häufige Muster im Trainingskorpus geprägt sind als durch die tatsächlichen Eigenschaften des analysierten Datensatzes \parencite{wang_generalization_2024}. \textcite{bender_dangers_2021} bezeichnen \glspl{llm} in diesem Zusammenhang als stochastische Papageien: Modelle, die sprachliche Formen reproduzieren, ohne notwendigerweise die zugrundeliegenden Inhalte zu verstehen.

Die operative Steuerung dieses Modellverhaltens erfolgt zur Laufzeit maßgeblich durch den Eingabe-Prompt. Durch In-Context Learning kann ein Modell allein auf Basis von Beispielen oder Anweisungen im Prompt auf neue Aufgaben ausgerichtet werden, ohne dass eine Anpassung der Modellgewichte erforderlich ist \parencite{brown_language_2020}. Eine für analytische Aufgaben besonders relevante Strategie ist Chain-of-Thought-Prompting, bei dem das Modell angewiesen wird, Zwischenschritte explizit auszuformulieren, bevor es eine Antwort generiert \parencite{wei_chain--thought_2023}. Im Kontext der \gls{eda} erweitert sich diese Prompt-Ebene zu einer iterativen Mensch-Modell-Interaktion. Da jede Modellantwort in den nachfolgenden Prompt einfließt, entstehen mehrstufige Analyseprozesse, in denen frühe Ungenauigkeiten in späteren Schritten fortgeführt oder verstärkt werden können \parencite{langer_towards_2019}.

Ab einer bestimmten Modellgröße entstehen Fähigkeiten, die bei kleineren Modellen nicht beobachtet werden und nicht linear aus dem Trainingsumfang ableitbar sind \parencite{wei_emergent_2022}. \todo{inwieweit dieses methodische Vorwissen für die Datenanalyse trägt, wird in dieser Arbeit untersucht. "Dies legt nahe, dass größere
Modelle implizit methodisches Wissen über statistische Verfahren erworben
haben können, ohne hierauf explizit trainiert worden zu sein."} Gleichzeitig ist bei der Verarbeitung strukturierter Daten schwer zu unterscheiden, ob Ausgaben auf echtem Generalisierungsvermögen oder auf der Reproduktion gelernter Muster beruhen \parencite{wang_generalization_2024}.
 
\section{LLMs in der Datenanalyse}
\label{sec:llm-analyse}

Die zuvor beschriebenen Fähigkeiten von \glspl{llm} werden auch für Aufgaben der Datenanalyse genutzt \parencite{cheng_gpt_2023, li_tapilot-crossing_2024}. Grundlage hierfür ist ihre Fähigkeit, natürlichsprachliche Anfragen in ausführbaren Programmcode zu überführen \parencite{chen_evaluating_2021}. Damit verschiebt sich der Untersuchungsgegenstand von den sprachlichen Fähigkeiten der Modelle hin zu ihrer Anwendung auf strukturierten Daten.

Einen umfassenden Ansatz in dieser Richtung verfolgen \textcite{cheng_gpt_2023}, die GPT-4 als eigenständigen Datenanalysten einsetzen und den Arbeitsablauf in Datenextraktion, Visualisierung und Analyse gliedern. Ihr Ansatz folgt dem dreistufigen Schema aus Code-Generierung, Code-Ausführung und anschließender Analysegenerierung, das bei Bedarf um externes Wissen ergänzt wird. In ihren Experimenten erreicht das Modell eine mit erfahrenen menschlichen Analystinnen und Analysten vergleichbare Qualität und übertrifft das Niveau von Berufseinsteigern, bei zugleich geringeren Kosten und kürzerer Bearbeitungszeit \parencite{cheng_gpt_2023}.

Während dieser Ansatz von einer einmaligen, eindeutig formulierten Anfrage ausgeht, argumentieren \textcite{li_tapilot-crossing_2024}, dass reale Analyseprozesse interaktiv verlaufen, da Nutzerintentionen häufig uneindeutig sind und Analysestrategien an Zwischenergebnisse angepasst werden. Sie beschreiben interaktive Datenanalyse als eine Abfolge miteinander verknüpfter Nutzer-Agent-Interaktionen und unterscheiden verschiedene Agentenaktionen, darunter das Stellen von Rückfragen, das Treffen begründeter Annahmen sowie die Korrektur und Aktualisierung von Code. Zur Verbesserung dieser Agenten schlagen sie eine Reflexionsstrategie vor, die aus erfolgreichen Interaktionsverläufen lernt und die relative Leistung um bis zu 44,5\% steigert. Dabei werden Reasoning-Verfahren wie Chain-of-Thought \parencite{wei_chain--thought_2023} eingesetzt. Zugleich beobachten sie, dass ein zu starkes Verlassen auf die Interaktionshistorie das eigenständige Schließen aus dem aktuellen Kontext beeinträchtigen kann, sodass zwischen der Nutzung früherer Interaktionsverläufe und der Reaktion auf die unmittelbare Eingabe abzuwägen ist \parencite{li_tapilot-crossing_2024}.

Über solche Systeme hinaus wurden Benchmarks vorgeschlagen, die \gls{llm}-gestützte Datenanalyse systematisch bewerten, etwa InfiAgent-DABench \parencite{hu_infiagent-dabench_2024}. \todo{Bei Bedarf mehr Quellen erängzen} 

Über die reine Leistungsmessung hinaus rückt dabei die für diese Arbeit zentrale Unterscheidung zwischen datenbasierten und wissensbasierten Ausgaben in den Blick. \textcite{li_tapilot-crossing_2024} adressieren sie explizit, indem sie mit nutzerdefinierten Bibliotheken prüfen, ob ein Modell die bereitgestellten Funktionen tatsächlich verarbeitet oder auf memorierte Standardsyntax zurückgreift. Diese Spannung bildet einen zentralen Ausgangspunkt der vorliegenden Arbeit.

Speziell für die \gls{eda} stellt der Einsatz von \glspl{llm} eine jüngere Entwicklung dar, die klassische Analysemethoden nicht ersetzt, sondern um generative Fähigkeiten erweitert \parencite{zhecheva_exploratory_2024}. Empirische Untersuchungen zeigen, dass die Qualität der \gls{llm}-Ausgaben stark von der Struktur und dem Informationsgehalt des Eingabe-Prompts abhängt: Bei klar strukturierten Prompts, die den zugrundeliegenden Datensatz enthalten, stimmen die \gls{llm}-Ausgaben weitgehend mit unabhängig berechneten Referenzwerten überein. Aufgaben, die ein semantisches Verständnis der Datendomäne erfordern, werden hingegen deutlich schlechter bewältigt, was auf Grenzen der faktischen Verlässlichkeit bei Analyseaufgaben hinweist \parencite{zhecheva_exploratory_2024}.

Die dargestellten Arbeiten zeigen Anwendungen \glspl{llm}-gestützter Datenanalyse, verweisen jedoch auf Einschränkungen wie fehlerhafte Berechnungen \parencite{cheng_gpt_2023} oder unbegründete Annahmen \parencite{li_tapilot-crossing_2024}. Diese Zuverlässigkeitsaspekte werden im folgenden Abschnitt vertieft.

\section{Zuverlässigkeit und Grenzen LLM-gestützter Analyse}

Bevor die spezifischen Fehlerquellen \gls{llm}-gestützter Analyse betrachtet werden, ist zu klären, worin die Zuverlässigkeit eines analytischen Systems besteht. \textcite{bornschlegl_towards_2024} überträgt die Dimensionen zwischenmenschlichen Vertrauens auf die Vertrauenswürdigkeit technischer Systeme und benennt für deren Ausgaben insbesondere Reproduzierbarkeit und Validität. Reproduzierbarkeit bezeichnet die Fähigkeit, mit denselben Eingabedaten, Rechenschritten, Methoden und demselben Code konsistente Ergebnisse zu erzielen. Validität beschreibt den Grad, in dem die erzeugten Daten oder Ergebnisse korrekt sind. Übertragen auf die \gls{llm}-gestützte explorative Datenanalyse beschreibt Reproduzierbarkeit demnach die Stabilität der Ergebnisse bei gleicher Eingabe und Validität ihre Übereinstimmung mit den tatsächlichen Eigenschaften der Daten \parencite{bornschlegl_towards_2024}. Inwieweit diese beiden Anforderungen durch charakteristische Eigenschaften von \glspl{llm} beeinträchtigt werden können, ist Gegenstand der folgenden Betrachtungen.

\textcite{ji_ziwei_survey_2022} weisen darauf hin, dass moderne Deep-Learning-Modelle zwar sprachlich kohärente Ausgaben erzeugen, jedoch häufig Inhalte produzieren, die weder mit der grundlegenden Eingabequelle übereinstimmen noch faktisch korrekt sind. Da solche als Halluzination bezeichnete Ausgaben definitionsgemäß von den tatsächlichen Fakten oder der bereitgestellten Eingabe abweichen, betreffen sie unmittelbar die zuvor eingeführte Validität. Solche Angaben lassen sich in zwei Typen unterscheiden. Faktische Halluzination bezeichnet die Diskrepanz zwischen generiertem Inhalt und verifizierbaren Fakten, während treuebezogene Halluzinationen die Abweichung von der Nutzeranweisung oder dem bereitgestellten Eingabekontext sowie eine fehlende Selbstkonsistenz umfassen \parencite{bai_hallucination_2025}. Für die vorliegende Arbeit ist vor allem der zweite Typ relevant, da der analysierte Datensatz den Eingabekontext bildet, von dem eine treuebezogene Halluzination abweicht. Erschwerend kommt hinzu, dass halluzinierte Ausgaben aufgrund ihrer sprachlichen Flüssigkeit und scheinbaren Kontextverankerung kaum von faktisch korrekten Inhalten zu unterscheiden sind \parencite{ji_ziwei_survey_2022}.

Dass Halluzinationen messbar und modellabhängig sind, zeigt eine Untersuchung zu fabrizierten Literaturangaben. \textcite{agrawal_language_2024} bestimmen für die Generierung wissenschaftlicher Referenzen modellspezifische Halluzinationsraten, die zwischen 47 Prozent für GPT-4 und über 76 Prozent für kleinere Modelle liegen, wobei leistungsfähigere Modelle durchgängig geringere Raten aufweisen. Zugleich betonen sie, dass Halluzination kein rein binäres Merkmal darstellt, da sich etwa eine generierte Referenz mit Teilen eines existierenden Titels nicht eindeutig klassifizieren lässt \parencite{agrawal_language_2024}. Ein anschauliches Beispiel aus der explorativen Datenanalyse ist der von GPT-4 erzeugte, nicht existente Titel \enquote{Exploratory Data Analysis and the Role of Visualization}, dem das Modell zudem den Autor \textit{John W. Tukey} der klassischen Arbeit zur explorativen Datenanalyse zuschrieb \parencite{agrawal_language_2024}.

Über die empirische Beobachtung hinaus wird die Unvermeidbarkeit von Halluzination theoretisch begründet. \textcite{xu_hallucination_2025} zeigen mit lerntheoretischen Mitteln, dass \glspl{llm} nicht alle berechenbaren Funktionen erlernen können und als allgemeine Problemlöser daher zwangsläufig halluzinieren. Für die Analysepraxis bedeutsam ist ihre Schlussfolgerung, dass sich Halluzination nicht allein durch veränderte Prompts beseitigen lässt und die Wirksamkeit des in Abschnitt \ref{sec:llm-analyse} beschriebenen Chain-of-Thought-Promptings somit auf bestimmte Aufgabentypen begrenzt bleibt. Als wirksamerer Ansatzpunkt gilt die Anreicherung mit externem Wissen, etwa durch Wissensdatenbanken oder Werkzeuge, da dem Modell so Information über die Zielfunktion zugeführt wird, die über die reinen Trainingsdaten hinausgeht \parencite{xu_hallucination_2025}. In dieselbe Richtung weist die Untersuchung von \textcite{zheng_reliable_2024} zur Verlässlichkeit von \glspl{llm} als Wissensbasis. Ein faktentreues Modell gibt bekanntes Wissen korrekt wieder und stellt zu unbekanntem keine Behauptungen auf, wobei sich die Überschätzung des eigenen Wissens besonders bei numerischen und zeitbezogenen Fragen zeigt. Konsistenz allein ist dabei kein hinreichendes Gütekriterium, da auch fehlerhafte Antworten konsistent wiederholt werden können \parencite{zheng_reliable_2024}.

Besondere Bedeutung erhält in diesem Zusammenhang die in Abschnitt \ref{sec:llm} beschriebene rein parametrische Wissensrepräsentation. Ein Modell unterscheidet nicht zwischen erlerntem Vorwissen und den im Prompt bereitgestellten Daten \parencite{wang_generalization_2024} und kann einen bekannten Datensatz daher vermeintlich korrekt beschreiben, ohne die konkret bereitgestellten Werte praktisch zu verarbeiten. Dieser Konflikt zwischen datenbasierten und wissensbasierten Ausgaben markiert eine zentrale Grenze der Verlässlichkeit \gls{llm}-gestützter Datenanalyse.

Zusammenfassend zeigen die dargestellten Eigenschaften, dass \glspl{llm} die benannten Anforderungen an Reproduzierbarkeit und Validität nur eingeschränkt erfüllen. Halluzinationen gefährden die Validität, indem sie inhaltlich nicht gedeckte Ausgaben erzeugen und ihre begründete Unvermeidbarkeit schließt eine vollständige Beseitigung aus \parencite{xu_hallucination_2025}. Hinzu kommt, dass ein Modell aufgrund seiner rein parametrischen Wissensrepräsentation dazu neigen kann, Ausgaben aus vortrainiertem Wissen statt aus den bereitgestellten Daten zu generieren. Für die \gls{llm}-gestützte explorative Datenanalyse folgt daraus, dass die Übereinstimmung der Modellausgaben mit dem tatsächlichen Dateninhalt nicht vorausgesetzt, sondern eigens geprüft werden muss. Im Kontext dieser Arbeit wird Datenkonsistenz als die Übereinstimmung zwischen einer generierten Aussage und dem zugrunde liegenden Datensatz verstanden. Sie stellt eine wesentliche Voraussetzung für die Beurteilung der Verlässlichkeit generierter Inhalte dar. 